\clearpage
\section{선별 문제 풀이 I}
\label{sec:norm-trace-solutions-a}

\subsection*{문제 \thechapter.1: 이차확장의 행렬 계산}

\begin{solution}
$u=3+2\sqrt5$라 하자. 기저 $(1,\sqrt5)$에 대해
\[
  u\cdot1=3+2\sqrt5,
\]
\[
  u\cdot\sqrt5=10+3\sqrt5.
\]
따라서
\[
  M_u=
  \begin{pmatrix}
    3&10\\
    2&3
  \end{pmatrix}.
\]
특성다항식은
\[
\begin{aligned}
  \chi_u(X)
  &=\det
  \begin{pmatrix}
    X-3&-10\\
    -2&X-3
  \end{pmatrix}\\
  &=(X-3)^2-20\\
  &=X^2-6X-11.
\end{aligned}
\]
그러므로
\[
  \boxed{\Tr_{L/\Q}(u)=6},
  \qquad
  \boxed{\Nm_{L/\Q}(u)=-11}.
\]
비자명한 켤레는 $3-2\sqrt5$이므로
\[
  (3+2\sqrt5)+(3-2\sqrt5)=6,
\]
\[
  (3+2\sqrt5)(3-2\sqrt5)=9-20=-11
\]
로 검산된다.
\end{solution}

\subsection*{문제 \thechapter.3: 순수 삼차확장}

\begin{solution}
$X^3-2$는 아이젠슈타인 판정으로 $\Q$ 위에서 기약이다.
따라서 $[L:\Q]=3$이고 $\alpha$의 최소다항식은 $X^3-2$이다.
\[
  \boxed{\Tr(\alpha)=0,\qquad \Nm(\alpha)=2.}
\]

$\alpha^2$는 $(\alpha^2)^3=4$를 만족한다. 삼차다항식 $X^3-4$는
유리근을 갖지 않으므로 $\Q$ 위에서 기약이다. 따라서 최소다항식은 $X^3-4$이다.
\[
  \boxed{\Tr(\alpha^2)=0,\qquad \Nm(\alpha^2)=4.}
\]

$\beta=1+\alpha$라 두면 $\alpha=\beta-1$이므로
\[
  (\beta-1)^3-2=0.
\]
따라서 $\beta$의 최소다항식은
\[
  X^3-3X^2+3X-3.
\]
차수가 이미 $3$이므로
\[
  \boxed{\Tr(1+\alpha)=3,\qquad \Nm(1+\alpha)=3.}
\]
\end{solution}

\subsection*{문제 \thechapter.4: 이중 이차확장의 탑 계산}

\begin{solution}
$E=\Q(\sqrt2)$라 하자. $[L:E]=2$이고 $L/E$의 비자명한 자동동형은
$\sqrt3$의 부호를 바꾸며 $\sqrt2$를 고정한다.

먼저 $\sqrt2\in E$이므로
\[
  \Tr_{L/E}(\sqrt2)=2\sqrt2,
  \qquad
  \Nm_{L/E}(\sqrt2)=2.
\]
탑 공식을 적용하면
\[
  \Tr_{L/\Q}(\sqrt2)
  =\Tr_{E/\Q}(2\sqrt2)=0,
\]
\[
  \Nm_{L/\Q}(\sqrt2)
  =\Nm_{E/\Q}(2)=2^2=4.
\]

이제 $u=\sqrt2+\sqrt3$에 대해
\[
  \Tr_{L/E}(u)
  =(\sqrt2+\sqrt3)+(\sqrt2-\sqrt3)
  =2\sqrt2,
\]
\[
  \Nm_{L/E}(u)
  =(\sqrt2+\sqrt3)(\sqrt2-\sqrt3)
  =-1.
\]
따라서
\[
  \boxed{\Tr_{L/\Q}(u)=0},
  \qquad
  \boxed{\Nm_{L/\Q}(u)=1}.
\]

직접 제곱하면
\[
  u^2=5+2\sqrt6,\qquad
  (u^2-5)^2=24,
\]
이므로
\[
  u^4-10u^2+1=0.
\]
$u$가 $L$을 생성하므로 최소다항식은 $X^4-10X^2+1$이고,
그 계수에서도 같은 대각합과 노름을 읽을 수 있다.
\end{solution}

\subsection*{문제 \thechapter.5: $\F_8/\F_2$의 대각합과 노름}

\begin{solution}
$\alpha^3=\alpha+1$이고
\[
  \alpha^4=\alpha^2+\alpha,\quad
  \alpha^5=\alpha^2+\alpha+1,\quad
  \alpha^6=\alpha^2+1,\quad
  \alpha^7=1.
\]
대각합은
\[
  \Tr(x)=x+x^2+x^4
\]
이다. $\F_2$-선형성을 이용하면
\[
  \Tr(1)=1,\qquad
  \Tr(\alpha)=0,\qquad
  \Tr(\alpha^2)=0.
\]
따라서 다음 표를 얻는다.
\[
\begin{array}{c|cccccccc}
x&0&1&\alpha&\alpha^2&
1+\alpha&1+\alpha^2&
\alpha+\alpha^2&1+\alpha+\alpha^2\\
\hline
\Tr(x)&0&1&0&0&1&1&0&1
\end{array}
\]
그러므로
\[
  \boxed{
  \ker\Tr
  =\{0,\alpha,\alpha^2,\alpha+\alpha^2\}.}
\]

한편 $\F_8^\times$의 위수는 $7$이고
\[
  \Nm_{\F_8/\F_2}(x)=x^{1+2+4}=x^7.
\]
따라서 모든 $x\ne0$에 대해
\[
  \boxed{\Nm(x)=1},
\]
그리고 $\Nm(0)=0$이다.
\end{solution}

\subsection*{문제 \thechapter.8: 이차확장의 쌍대기저}

\begin{solution}
기저 $(1,\sqrt d)$에 대해
\[
  \Tr(1)=2,\qquad
  \Tr(\sqrt d)=0,\qquad
  \Tr(d)=2d.
\]
따라서 그람행렬은
\[
  G=
  \begin{pmatrix}
    2&0\\
    0&2d
  \end{pmatrix},
\qquad
  \det G=4d.
\]
$d\ne0$이고 $\charac K\ne2$이므로 판별식은 영이 아니다.

$y_1=a+b\sqrt d$가
\[
  \Tr(y_1)=1,\qquad
  \Tr(\sqrt d\,y_1)=0
\]
을 만족해야 한다. 첫 식은 $2a=1$, 둘째 식은 $2bd=0$이므로
\[
  y_1=\frac12.
\]
마찬가지로 $y_2=a+b\sqrt d$에 대해
\[
  \Tr(y_2)=0,\qquad
  \Tr(\sqrt d\,y_2)=1
\]
에서 $a=0$, $2bd=1$이므로
\[
  y_2=\frac{\sqrt d}{2d}.
\]
따라서
\[
  \boxed{
  \left(\frac12,\frac{\sqrt d}{2d}\right)}
\]
가 쌍대기저이다.
\end{solution}

\subsection*{문제 \thechapter.11: 아르틴의 선형독립 정리}

\begin{solution}
서로 다른 캐릭터들 사이에 비자명한 선형관계가 있다고 가정하고,
그 가운데 항의 수 $r$이 최소인 것을
\[
  c_1\chi_1+\cdots+c_r\chi_r=0
\]
로 택한다. 최소성에서 모든 $c_i\ne0$이다.

$\chi_1\ne\chi_2$이므로 어떤 $g\in G$에 대해
$\chi_1(g)\ne\chi_2(g)$이다. 원래 관계를 $gh$에 적용하면
\[
  \sum_{i=1}^r c_i\chi_i(g)\chi_i(h)=0
\]
이고, 이는 함수관계
\[
  \sum_{i=1}^r c_i\chi_i(g)\chi_i=0
\]
를 뜻한다. 원래 관계에 $\chi_1(g)$를 곱한 것에서 이 관계를 빼면
\[
  \sum_{i=2}^r
  c_i\bigl(\chi_1(g)-\chi_i(g)\bigr)\chi_i=0.
\]
$\chi_2$의 계수
\[
  c_2\bigl(\chi_1(g)-\chi_2(g)\bigr)
\]
는 영이 아니므로 길이 $r-1$의 비자명한 관계를 얻는다.
이는 $r$의 최소성에 모순이다. 따라서 서로 다른 캐릭터는 선형독립이다.
\end{solution}

\subsection*{문제 \thechapter.14: 최소다항식에서 노름과 대각합 읽기}

\begin{solution}
모든 $h\in K[X]$에 대해
\[
  h(m_\alpha)=m_{h(\alpha)}
\]
이다. 따라서 $h(\alpha)=0$이면 $h(m_\alpha)=0$이다.
반대로 $h(m_\alpha)=0$이면 이를 $1$에 적용하여
\[
  h(\alpha)=h(m_\alpha)(1)=0
\]
을 얻는다. 그러므로 $\alpha$와 $m_\alpha$의 최소다항식은 같다.

$L=K(\alpha)$이고 차수가 $n$이므로 이 최소다항식의 차수는 $n$이다.
한편 $m_\alpha$의 특성다항식도 $n$차이고 최소다항식으로 나누어지므로
두 다항식은 같다.
\[
  \chi_{\alpha,L/K}(X)
  =X^n+c_1X^{n-1}+\cdots+c_n.
\]
특성다항식의 일반형
\[
  X^n-\Tr(\alpha)X^{n-1}
  +\cdots+(-1)^n\Nm(\alpha)
\]
과 비교하면
\[
  \boxed{\Tr_{L/K}(\alpha)=-c_1},
  \qquad
  \boxed{\Nm_{L/K}(\alpha)=(-1)^nc_n}.
\]
\end{solution}

\subsection*{문제 \thechapter.16: 체매장에 의한 공식}

\begin{solution}
먼저 $L=K(a)$라 하자. 서로 다른 $K$-매장의 수를 $r$,
비분리차수를 $q$라 하면 $a$의 최소다항식은
\[
  f_{a,K}(X)
  =\prod_{i=1}^r(X-\sigma_i(a))^q.
\]
이는 각 서로 다른 켤레가 정확히 비분리차수 $q$만큼의 중복도를
갖기 때문이다. $L=K(a)$이므로 이 최소다항식은 $a$의
$L/K$-특성다항식과 같다. $X^{[L:K]-1}$의 계수와 상수항을 비교하면
\[
  \Tr_{L/K}(a)
  =q\sum_{i=1}^r\sigma_i(a),
\]
\[
  \Nm_{L/K}(a)
  =\left(\prod_{i=1}^r\sigma_i(a)\right)^q.
\]

이제 일반적인 $a\in L$에 대해 $E=K(a)$, $s=[L:E]$라 하자.
블록대각행렬 계산에서
\[
  \Tr_{L/K}(a)=s\Tr_{E/K}(a),
  \qquad
  \Nm_{L/K}(a)=\Nm_{E/K}(a)^s.
\]
각 $K$-매장 $E\to\overline K$는 정확히 $[L:E]_{\mathrm s}$개의
$K$-매장 $L\to\overline K$로 연장되고, 모든 연장은 $a$에서 같은 값을 갖는다.
또한
\[
  s=[L:E]_{\mathrm s}[L:E]_{\mathrm i}
\]
이고 비분리차수는 탑에서 곱셈적이다. 따라서 $E/K$에서 이미 얻은
합과 곱의 공식에 이 중복도를 반영하면 동일한 두 공식을 $L/K$에 대해 얻는다.
\end{solution}
